\section{谱作用}\srcpage{25}

谱作用在各种情形下给出 $D_{\qc}(\LocSys_{\Gd})$ 在 $D(\Bun_G)$ 上的一个作用。为得到它，可以更一般地追问：给出 \srcpage{26}$D_{\qc}(\LocSys_{\Gd})$ 的一个作用，究竟需要什么。在 Betti 情形，这有一个漂亮的答案，归功于 Nadler 与 Yun（2019）：

\begin{theorem}
设 $C\in\PrL_{\e}$。在 $C$ 上给出 $D_{\qc}(\LocSys^{\Betti}_{\Gd})$ 的一个作用，等价于给出幺半 $e$-线性范畴之间的正合函子
\[
\Rep(\Gd^I)\to\End(C)^{B\pi_1^I},
\]
且关于有限集 $I$ 函子式地成立。这里，右端表示幺半范畴 $\End(C)$ 中的 $\pi_1^I$-等变对象，而 $\End(C)$ 是 $C$ 的 $e$-线性、保余极限自函子构成的范畴。
\end{theorem}

这里，我们用 $\Rep(\Gd)$ 表示 $\Gd$ 的有限维表示构成的小 Abel 范畴。当从它出发作映射时，我们隐含地用到：对可表现稳定的 $C$，正合函子 $\Rep(\Gd)\to C$ 与保余极限函子 $D_{\qc}(\ast/\Gd)\to C$ 是同一回事。

\begin{remark}
定理 5.1 有一个更初等的类似物：它描述的不是 $\LocSys^{\Betti}_{\Gd}$ 上的范畴层，而是模层，即拟凝聚层。这可以用同样的方式证明，结论是：对象 $M\in D_{\qc}(\LocSys^{\Betti}_{\Gd})$ 等价于函子
\[
\Rep(\Gd^I)\to D(\ast/\pi_1^I)=D(e)^{B\pi_1^I}
\]
的数据，且关于 $I$ 函子式。这把 $M$ 送到对应
\[
\begin{tikzcd}
\LocSys^{\Betti}_{\Gd}\times\ast/\pi_1^I \arrow[r] \arrow[d] & \ast/\Gd^I \\
\ast/\pi_1^I
\end{tikzcd}
\]
上以 $M$ 为核的函子。这里，上方的箭头伴随于
\[
\LocSys^{\Betti}_{\Gd}=\Map(\ast/\pi_1,\ast/\Gd)\to\Map(\ast/\pi_1^I,\ast/\Gd^I).
\]
\end{remark}

\begin{proof}[证明概要]
我们给出一个论证，所用的是 Gestalt（格式塔）的语言，取自 Scholze（2026）关于与 Stefanich 合作工作的课程。$e$ 上的一个 Gestalt $X$ 由一列 $(\cO(X)_0,\cO(X)_1,\dots)$ 给出。其中 $\cO(X)_0$ 是 $D(e)$ 中的交换环，即「$X$ 上的函数环」。$\cO(X)_1$ 是幺半 $e$-线性 $(\infty,1)$-范畴，其单位的自同态给出 $\cO(X)_0$，即「$X$ 上的拟凝聚层范畴」。$\cO(X)_2$ 是幺半 $e$-线性 $(\infty,2)$-范畴，其单位的自同态给出 $\cO(X)_1$，即「$X$ 上的线性范畴层构成的 2-范畴」。……因此，Gestalten 旨在捕捉所有范畴层级上的高阶代数。出人意料的结论是：它们本质上构成一个 $(\infty,1)$-拓扑斯，从而可以像几何对象一样被思考与处理。因此，Gestalten 的语言把高阶范畴中困难的代数，翻译为更易驾驭的几何。这与概形的语言把困难的交换代数翻译为更易驾驭的几何，方式相同。

\srcpage{27}在 $e$ 上的 Gestalten 中，考虑映射 Gestalt
\[
\Map(\ast/\pi_1,\ast/\Gd)=\Hom(\pi_1,\Gd)/\Gd.
\]
这实际上就是 $\LocSys^{\Betti}_{\Gd}$，经由从代数叠到 Gestalten 的自然函子而视为一个 Gestalt。另一方面，映射 Gestalt 有一个一般公式。即：若 $X$ 是 $n$-suave 的、$Y$ 是 $n$-affine 的，则映射 Gestalt $\Map(X,Y)$ 自身是 $n$-affine 的，从而对应某个对称幺半可表现 $(\infty,n)$-范畴
\[
\cO(\Map(X,Y))_n\in\operatorname{CAlg}(\mathrm{nPr}_e);
\]
这里，$\mathrm{nPr}_e$ 表示可表现 $e$-线性 $(\infty,n)$-范畴构成的对称幺半范畴。此时 $\cO(\Map(X,Y))_n$ 实际上是万有的 $C\in\operatorname{CAlg}(\mathrm{nPr}_e)$，它配有对称幺半函子
\[
\cO(Y)_n\to\Hom(\cO(X)^!_n,C).
\]
这可由 $\Map(X,Y)$ 的定义以及 suave 映射的基变换公式推出（Scholze，2026，命题 6.18），后者涉及 $\cO(X)_n$ 的一个前对偶 $\cO(X)^!_n$。若忘掉 $C$ 上的交换代数结构，则万有的 $C$ 会是 $\cO(Y)_n\otimes\cO(X)^!_n$。回忆：Lurie 高阶代数中交换代数的概念，以有限集的操作子来编码。因此，万有的对称幺半 $C$ 可写成如下函子的余端：它把任一有限集 $I$ 送到
\[
\cO(Y)^{\otimes I}_n\otimes(\cO(X)^!_n)^{\otimes I}.
\]

在我们的情形，$X=\ast/\pi_1$ 是 $0$-étale 的，从而对任意 $n\ge 0$ 都是 $n$-suave 的；而 $Y=\ast/\Gd$ 是 $1$-affine 的。因此，该映射 Gestalt 是 $1$-affine 的。事实上，我们已知它由代数叠 $\LocSys^{\Betti}_{\Gd}=\Hom(\pi_1,\Gd)/\Gd$ 给出，所以
\[
\Map(\ast/\pi_1,\ast/\Gd)=\mathrm{Gest}(D_{\qc}(\LocSys^{\Betti}_{\Gd})).
\]
特别地，由于它是 $1$-affine 的，有
\[
\cO(\Map(\ast/\pi_1,\ast/\Gd))_2=\operatorname{Mod}_{D_{\qc}(\LocSys^{\Betti}_{\Gd})}(\mathrm{1Pr}_e),
\]
而这正是我们想理解的。按上述一般描述，它可描述为如下函子的余端：
\[
I\mapsto\cO(\ast/\Gd)^{\otimes I}_2\otimes(\cO(\ast/\pi_1)^!_2)^{\otimes I}.
\]
这里，$\cO(\ast/\Gd)^{\otimes I}_2=\operatorname{Mod}_{\Rep(\Gd)^{\otimes I}}(\mathrm{1Pr}_e)$，而加上第二个张量因子后，它就成为其中的 $\pi_1^I$-等变对象。拆解开来即得定理。
\end{proof}

\srcpage{28}定理 5.1 的数据几乎就是 Hecke 作用给我们的东西！更确切地说，对任一有限集 $I$ 与 $X^I$ 中任意的 $I$ 个点 $(x_i)_{i\in I}$，我们有 Hecke 对应
\[
\begin{tikzcd}
& \Hck_{G,(x_i)_i} \arrow[ld] \arrow[rd] & \\
\Bun_G & & \Bun_G.
\end{tikzcd}
\]
利用 $I$ 个点处几何 Satake 等价的一个版本，我们得到 $e$-线性幺半函子
\[
\Rep(\Gd^I)\to\Perv(\Hck_{G,(x_i)_i}),
\]
它关于目标上的卷积结构是幺半的。另一方面，以 Hecke 叠上的反常层作核，我们得到 $e$-线性幺半函子
\[
\Perv(\Hck_{G,(x_i)_i})\to\End(D(\Bun_G)).
\]
复合，便得到 $e$-线性幺半函子
\[
\Rep(\Gd^I)\to\End(D(\Bun_G)).
\]
但现在我们可以让点 $(x_i)_{i\in I}\in X^I$ 移动。若上述函子在点变动时适当地局部常值，就会给出所要的映射
\[
\Rep(\Gd^I)\to\End(D(\Bun_G))^{B\pi_1^I},
\]
即一个 $\pi_1^I=\pi_1(X^I)$-作用。

更确切地说，Hecke 作用实际上给出函子
\[
\Rep(\Gd^I)\times D(\Bun_G)\to D(\Bun_G\times X^I).
\]
在目标上，我们有沿 $X^I$-方向局部常值的层构成的全子范畴。它们等价地就是
\[
D(\Bun_G)^{B\pi_1^I}\cong D(\Bun_G\times\ast/\pi_1^I)\to D(\Bun_G\times X^I)
\]
的本质像中的层，这里用的是沿 $X^I\to\ast/\pi_1^I$ 的拉回。

正如我们在定理 3.4 处所讨论的，这对 $\Bun_G$ 上的所有层并不成立，但对
\[
D_{\Nilp}(\Bun_G)\subset D(\Bun_G)
\]
确实成立。在 étale 情形，情况完全平行；定理 5.1 有直接的类比成立，参见 Arinkin、Gaitsgory、Kazhdan、Raskin、Rozenblyum 与 Varshavsky（2020b，主定理 8.1.4）。

\begin{corollary}[Nadler 与 Yun，2019；Arinkin、Gaitsgory、Kazhdan、Raskin、Rozenblyum 与 Varshavsky，2020b]
在 Betti 或 étale 情形，存在 $D_{\qc}(\LocSys_{\Gd})$ 在 $D_{\Nilp}(\Bun_G)$ 上的谱作用（其中一个来源见主定理 14.3.2）。
\end{corollary}

\begin{remark}
\srcpage{29}正如我们在猜想 1.5 之后所讨论的，这给出几何 Langlands 的「自守到 Galois」方向。Lafforgue（2018）用非常类似的想法，建立了函数域上算术 Langlands 的「自守到 Galois」方向。更确切地说，本注中仍设 $X$ 是 $\Fq$ 上的曲线。\footnote{与几何 Langlands 纲领的当前状况不同，对下面的讨论我们甚至可以允许分歧，正如 V. Lafforgue 所做的那样。}于是我们可以把 $\cA_c(\Bun_G(\Fq))$ 提升为（算术的）$\LocSys_{\Gd}$ 上的模，另参见 Lafforgue 与 Zhu（2019）。由注 5.2 的一个版本，这等价于给出函子
\[
\Rep(\Gd^I)\to D(\Ql)^{BW^I}
\]
其中 $W$ 表示有限域上该曲线的 Weil 群，即几何基本群被 $\Z$ 的半直积。这恰好由取 shtuka 空间的上同调这一函子给出！下面给出 shtuka 模空间的一个稍显非标准的表述，它却简化了关于「部分 Frobenius」的讨论。回忆：直观上，人们把映射
\[
X\to\Spec(\Fq)
\]
想成一个纤维化于圆周之上的三维流形。有一个简便的办法把它变成真正的数学。我们令
\[
\Spec(\Fq)^{\Weil}=(\Spec(\overline{\F}_q)\times\R_{>0})/(\Frob_{\Fq}\times q)^{\Z},
\]
其中 $\R_{>0}$ 作为凝聚态集进入（即作为 pro-étale 代数空间，把它写成局部投射有限集的商）。于是 $\Spec(\Fq)^{\Weil}$ 的底层拓扑空间是一个圆周 $S^1=\R_{>0}/q^{\Z}$，所有剩余域都是 $\overline{\F}_q$；但绕圆周走一圈，就作用以 Frobenius 自同构。类似地，可以定义
\[
X^{\Weil}=(X_{\overline{\F}_q}\times\R_{>0})/(\Frob_{\overline{\F}_q}\times q)^{\Z}\to\Spec(\Fq)^{\Weil},
\]
其底层拓扑空间现在纤维化于 $S^1$ 之上，所有纤维都给出 $X_{\overline{\F}_q}$；绕圆周走一圈，就作用以 $\overline{\F}_q$ 的 Frobenius 自同构。更一般地，对 $\overline{\F}_q$ 上任一完满概形 $S$，可类似地定义
\[
X^{\Weil}_S=(X_S\times\R_{>0})/(\Frob_S\times q)^{\Z},
\]
它几乎、但并非完全是 $X^{\Weil}$ 的基变换（因为 $X^{\Weil}_S$ 事实上并不映射到 $S$）。我们还需要 Frobenius 改作用在 $X$ 上的版本，
\[
X^W=(X_{\overline{\F}_q}\times\R_{>0})/(\Frob_X\times q^{-1})^{\Z},
\]
（这里我们隐含地过渡到 $X$ 的完满化）。注意 $\pi_1(X^W)=W$ 就是 $X$ 的 Weil 群。
\end{remark}

\begin{definition}
带有 $I$ 条 leg 的 shtuka 模空间，是模空间
\[
f^I\colon\Sht^I\to(X^W)^I
\]
\srcpage{30}它把 $\overline{\F}_q$ 上任一完满概形 $S$ 送到 $I$ 个映射 $x_i\colon S\to X^W$，以及一个 $G$-丛\footnote{有一个稍显混乱、但本质上是自动的翻译，把 $x_i$ 的图像变成 $X^{\Weil}_S$ 的闭子集。}，后者定义在 $X^{\Weil}_S{\setminus}\bigcup_i\{\Gamma_{x_i}\}$ 上。
\end{definition}

因此，这参数化 $I$ 重穿孔三维流形上的 $G$-丛。这里的重要之处在于：沿 $\R_{>0}$-方向，这些 $G$-丛必须是常值的 —— 对任一概形 $S$，$S\times[0,1]$ 上的向量丛等价于 $S$ 上的向量丛。然而，当你越过一个穿孔时，可能发生一次修改。把上述 shtuka 的定义与更经典的定义联系起来，是一件令人愉快的练习。上述定义能够把关于修改种类的各种 shtuka 变体，合并为单个对象。\footnote{提示：对 $S=\Spec(\Fq)$，可以通过观察 $S^1$ 上的纤维来理解 $X^{\Weil}{\setminus}\{x_i\}$ 上的向量丛。在大多数纤维上，这不过是整个 $X_{\overline{\F}_q}$ 上的向量丛。把 $S^1$ 上的点稍作变动，由上述同伦不变性，该向量丛保持常值。在某些纤维上，有些点缺失；在那些纤维上，可以在缺失点处发生修改。最后，若绕圆周走一圈，最终回到原来的向量丛，但被 Frobenius 扭了一下。}

对任一 $V\in\Rep(\Gd^I)$，几何 Satake 等价给出 $\Sht^I$ 上的层 $\Sat_V$，它由某个适当的局部 Hecke 叠拉回而来。由 Xue（2025）的一个定理，shtuka 模空间的上同调层是光滑的。此外，$\pi_1(X^W)=W$ 是 Weil 群，且 $X^W$ 是固有的，故关于 $\pi_1$ 的 Künneth 定理成立，且 $\pi_1((X^W)^I)=W^I$。因此，我们得到
\[
f^I_!\Sat_V\in D(\Ql)^{BW^I}\subset D((X^W)^I,\Ql),
\]
从而得到所要的函子
\[
\Rep(\Gd^I)\to D(\Ql)^{BW^I}.
\]

我期望存在 $\Spec(\Z)$ 上的 shtuka 理论 —— 建立在把 shtuka 视为穿孔三维流形上的向量丛这一模板上 —— 用它最终能在数域情形执行类似的论证。

\begin{remark}
在 Betti 情形，定理 5.1 有一个适用于整个 $D(\Bun_G)$ 的类似物。即：沿用几何 Langlands 纲领中使用的记号，\footnote{除了把 $\Rep(\Gd)$ 换成 $D_{\qc}(\ast/\Gd)$ 之外；我们用 $\Rep(\Gd)$ 表示那个小 Abel 范畴。「Ran」这个名字在这里指「Ran 空间」，即「$X$ 的一切有限子集构成的空间」。我通常避免使用 Ran 空间，因为它是性态有些不良的几何对象，这里只用该记号与文献衔接。}令
\[
D_{\qc}(\ast/\Gd)_{\Ran}
\]
表示万有的 $C\in\operatorname{CAlg}(\PrL_{\e})$，它配有对称幺半函子
\[
\Rep(\Gd)\to C\otimes D(X).
\]
\srcpage{31}那么 Hecke 作用总能产生 $D_{\qc}(\ast/\Gd)_{\Ran}$ 在 $D(\Bun_G)$ 上的一个作用。在 Scholze（2026）与 Stefanich 的工作所发展的 Gestalten 框架中，$C$ 对应映射 Gestalt
\[
\Map(\mathrm{Gest}(D(X)),\ast/\Gd),
\]
即从 $D(X)$ 对应的 Gestalt 出发、指向 $\Gd$ 的分类叠的映射。换言之，这就是「Betti $\Gd$-局部系统的 Gestalt」的自然概念，因此我们将记
\[
\LocSys^{\Betti,\mathrm{Gest}}_{\Gd}=\Map(\mathrm{Gest}(D(X)),\ast/\Gd)=\mathrm{Gest}(D_{\qc}(\ast/\Gd)_{\Ran}).
\]
于是整个 $D(\Bun_G)$ 承载 $D_{\qc}(\ast/\Gd)_{\Ran}$ 的一个作用。隐含地把代数叠与其相联系的 Gestalt 等同起来，有自然的单射
\[
\LocSys^{\Betti}_{\Gd}\hookrightarrow\LocSys^{\Betti,\mathrm{Gest}}_{\Gd}.
\]
遗憾的是，在几何 Langlands 下，$D_{\Nilp}(\Bun_G)$ 会本质等价于 $D_{\qc}(\LocSys_{\Gd})$；但对整个 $D(\Bun_G)$ 而言并非如此。它并不等价于 $D_{\qc}(\LocSys^{\Betti,\mathrm{Gest}}_{\Gd})=D_{\qc}(\ast/\Gd)_{\Ran}$。
\end{remark}

\begin{remark}
在 étale 情形同样有 $D_{\qc}(\ast/\Gd)_{\Ran}$ 在整个 $D(\Bun_G)$ 上的作用的类似物。最干净的陈述，是把 Gestalten 的形式化更认真地使用，这次一直走到 $(\infty,2)$-范畴。即：给定的六函子形式化 $D$ 给出一个 transmutation 函子
\[
\Sch^{\mathrm{ft}}_k\to\mathrm{Gest}/\mathrm{Gest}(D(e))\colon S\mapsto[S]_D,
\]
它把任一 $S$ 送到 $S$ 上核的可表现 $(\infty,2)$-范畴的 Gestalt，参见 Scholze（2026，第 9 讲）。由下降，这延拓到叠。于是还可以考虑映射 Gestalt
\[
\LocSys^{\et,\mathrm{Gest}}_{\Gd}=\Map_{[\ast]_D}([X]_D,\ast/\Gd).
\]
它在 $[\ast]_D$ 上是 $1$-affine 的，从而对应 $\Spec(k)$ 上核的可表现 $(\infty,2)$-范畴 $K_D$ 中的一个交换代数。后者又可以描述为余端；它或许可称为 $D_{\qc}(\ast/\Gd)_{\Ran}$。于是 $D(\Bun_G)$ 成为 Gestalt $\LocSys^{\et,\mathrm{Gest}}_{\Gd}$ 上的范畴层，即 $K_D$ 中 $D_{\qc}(\ast/\Gd)_{\Ran}$ 上的模。\footnote{这类推理即使在 motivic 六函子形式化中也行得通。}
\end{remark}

我们稍后将用到局部系统的 Gestalt 的如下性质。

\begin{proposition}
$\LocSys^{\et,\mathrm{Gest}}_{\Gd}\to[\ast]_D$ 的对角是 $0$-affine 的，而自然映射
\[
\LocSys^{\et}_{\Gd}\times_{\mathrm{Gest}(D(e))}[\ast]_D\to\LocSys^{\et,\mathrm{Gest}}_{\Gd}
\]
是 $1$-étale 且 $1$-affine 的单射。\footnote{注意：$\LocSys^{\et}_{\Gd}$ 这个受限的局部系统叠是非常大的连通分支的无交并，特别地不是可数个仿射概形的并。为此，应取 Scholze（2026）中所用的截断基数 $\kappa$ 大于连续统，好让 $1$-étale 性的断言在技术上成立。}
\end{proposition}

\begin{proof}
\srcpage{32}首先，我们证明对角是 $0$-affine 的。由于 $[X]_D\to[\ast]_D$ 是 $0$-suave 的（因为 $X$ 上同调光滑），而 $\ast/\Gd$ 的对角是 $0$-affine 的，只需证明下述更一般的断言：在任一基 Gestalt $B$ 之上，若有一个在 $B$ 上 $0$-suave 的 Gestalt $S$，且 $T'\to T$ 是 $0$-affine 的，则 $\Map_B(S,T')\to\Map_B(S,T)$ 是 $0$-affine 的。为此，固定一个映射 $f\colon S\to T$，并考虑这样的函子：它把 $B$ 上的任一 $B'$ 送到 $f|_{B'}\colon S\times_B B'\to T$ 在 $T'\to T$ 上的提升构成的 anima。由于 $T'\to T$ 是 $0$-affine 的，这就是 $\operatorname{CAlg}(\cO(T)_1)$ 中映射
\[
(\cO(T')/\cO(T))_0\to(\cO(S\times_B B')/\cO(T))_0
\]
的 anima。但 $(\cO(S\times_B B')/\cO(T))_0$ 是 $(\cO(S\times_B B')/\cO(S))_0$ 在遗忘函子 $\operatorname{CAlg}(\cO(S)_1)\to\operatorname{CAlg}(\cO(T)_1)$（即沿 $f$ 的前推）下的像。此外，由 suave 基变换，$(\cO(S\times_B B')/\cO(S))_0\in\operatorname{CAlg}(\cO(S)_1)$ 是 $(\cO(B')/\cO(B))_0\in\operatorname{CAlg}(\cO(B)_1)$ 沿 $0$-suave 映射 $S\to B$ 的拉回。因此可以看到，$(\cO(S\times_B B')/\cO(T))_0$ 只依赖于 $(\cO(B')/\cO(B))_0\in\operatorname{CAlg}(\cO(B)_1)$。从而 $B'$ 只通过其 $0$-affine 截断进入。于是，该模空间在 $B$ 上是 $0$-affine 的。

现在容易看出映射
\[
\LocSys^{\et}_{\Gd}\times_{\mathrm{Gest}(D(e))}[\ast]_D\to\LocSys^{\et,\mathrm{Gest}}_{\Gd}
\]
是单射。事实上，这意味着它的对角是同构；但该对角是 $0$-affine 的，而在此情形可以比较模空间的描述。由于源与目标都是 $1$-affine 的，该映射是 $1$-affine 的。又，源是代数叠，从而是 $1$-étale 的；而目标的对角是 $0$-affine 的，从而是 $1$-étale 的；合起来便推出该映射是 $1$-étale 的。
\end{proof}

在 D-模的情形，情况略有不同。一般形式化仍把 $D(\Bun_G)$ 变成
\[
\LocSys^{\dR,\mathrm{Gest}}_{\Gd}=\Map(X_{\dR},\ast/\Gd)=\mathrm{Gest}(D_{\qc}(\ast/\Gd)_{\Ran})
\]
上的范畴层，其中 $X_{\dR}=X/(X\subset X\times X)^{\wedge}$ 是 $X$ 的 Simpson de Rham 叠。同样有单射
\[
\LocSys^{\dR}_{\Gd}\hookrightarrow\LocSys^{\dR,\mathrm{Gest}}_{\Gd}.
\]
然而，这一次无需对 $D(\Bun_G)$ 作任何修改：

\begin{theorem}[Gaitsgory，2015；Gaitsgory 与 Raskin，2025b]
$\LocSys^{\dR,\mathrm{Gest}}_{\Gd}$ 上的范畴层 $D(\Bun_G)$ 支撑在子空间
\[
\LocSys^{\dR}_{\Gd}\hookrightarrow\LocSys^{\dR,\mathrm{Gest}}_{\Gd}
\]
上。等价地，$D_{\qc}(\ast/\Gd)_{\Ran}$ 在 $D(\Bun_G)$ 上的作用经过 Verdier 商
\[
D_{\qc}(\ast/\Gd)_{\Ran}\to D_{\qc}(\LocSys^{\dR}_{\Gd})
\]
分解。
\end{theorem}

\begin{remark}
与 Betti 或 étale 情形不同，容许 $D_{\qc}(\LocSys^{\dR}_{\Gd})$ 的一个作用，并不强制 Hecke 作用沿 $X^I$-方向局部常值。归根结底，这是因为函子
\[
\Rep(\Gd)\to D_{\qc}(\LocSys^{\dR}_{\Gd})\otimes D_{\qc}(X_{\dR})
\]
\srcpage{33}（它编码万有的 de Rham $\Gd$-局部系统）并不经过
\[
D_{\qc}(\LocSys^{\dR}_{\Gd})\otimes\mathrm{IndLisse}(X_{\dR})\subset D_{\qc}(\LocSys^{\dR}_{\Gd})\otimes D_{\qc}(X_{\dR})
\]
分解。而这又是因为：把 $k$-代数 $A$ 送到 $D(A)\otimes D_{\qc}(X_{\dR})$ 中 lisse 对象的函子，不与 $A$ 中的基变换交换。
\end{remark}

这个定理极其非平凡。如同 D-模情形中关于几何 Langlands 的许多定理，它用到了这一框架的一些特殊之处。值得注意的是，有一个归纳地生成 $D(\Bun_G)$ 的对象的办法，它经由两种程序：其一是熟悉的 Eisenstein 级数构造，其二是 Kac--Moody 局部化。这给出产生尖点对象的一般办法。在 Langlands 纲领中，这是实际产生尖点对象最通用的途径。归根结底，它是几何 Langlands 对应中困难的「Galois 到自守」方向的来源。就手头的定理而言，人们证明这两个来源生成整个 $D(\Bun_G)$，且两个来源都支撑在
\[
\LocSys^{\dR}_{\Gd}\hookrightarrow\LocSys^{\dR,\mathrm{Gest}}_{\Gd}
\]
上。

\begin{remark}
虽然（层范畴意义下的）$\LocSys^{\dR,\mathrm{Gest}}_{\Gd}$ 有干净的生成元与关系描述，但尚不知道如何通过施加更多关系来给出子空间 $\LocSys^{\dR}_{\Gd}$ 的描述。这使得直接检验 $D(\Bun_G)$ 是否支撑在该子空间上变得困难。
\end{remark}
